\documentclass[coursework]{SCWorks}
% Тип обучения (одно из значений):
%    bachelor   - бакалавриат (по умолчанию)
%    spec       - специальность
%    master     - магистратура
% Форма обучения (одно из значений):
%    och        - очное (по умолчанию)
%    zaoch      - заочное
% Тип работы (одно из значений):
%    coursework - курсовая работа (по умолчанию)
%    referat    - реферат
%  * otchet     - универсальный отчет
%  * nirjournal - журнал НИР
%  * digital    - итоговая работа для цифровой кафедры
%    diploma    - дипломная работа
%    pract      - отчет о научно-исследовательской работе
%    autoref    - автореферат выпускной работы
%    assignment - задание на выпускную квалификационную работу
%    review     - отзыв руководителя
%    critique   - рецензия на выпускную работу

% * Добавлены вручную. За вопросами к @mchernigin
\usepackage{preamble}

\begin{document}

% Кафедра (в родительном падеже)
\chair{математической кибернетики и компьютерных наук}

% Тема работы
\title{Разработка механизма редактирования учебных курсов для серверной части образовательной платформы <<MergeMinds>>}

% Курс
\course{2}

% Группа
\group{251}

% Факультет (в родительном падеже) (по умолчанию "факультета КНиИТ")
% \department{факультета КНиИТ}

% Специальность/направление код - наименование
% \napravlenie{02.03.02 "--- Фундаментальная информатика и информационные технологии}
% \napravlenie{02.03.01 "--- Математическое обеспечение и администрирование информационных систем}
% \napravlenie{09.03.01 "--- Информатика и вычислительная техника}
\napravlenie{09.03.04 "--- Программная инженерия}
% \napravlenie{10.05.01 "--- Компьютерная безопасность}

% Для студентки. Для работы студента следующая команда не нужна.
% \studentfemale

% Фамилия, имя, отчество в родительном падеже
\author{Тюменцева Радомира Александровича}

% Заведующий кафедрой 
\chtitle{доцент, к.\,ф.-м.\,н.}
\chname{С.\,В.\,Миронов}

% Руководитель ДПП ПП для цифровой кафедры (перекрывает заведующего кафедры)
% \chpretitle{
%     заведующий кафедрой математических основ информатики и олимпиадного\\
%     программирования на базе МАОУ <<Ф"=Т лицей №1>>
% }
% \chtitle{г. Саратов, к.\,ф.-м.\,н., доцент}
% \chname{Кондратова\, Ю.\,Н.}

% Научный руководитель (для реферата преподаватель проверяющий работу)
\satitle{ассистент} %должность, степень, звание
\saname{Н.\,С.\,Рыданов}

% Руководитель практики от организации (руководитель для цифровой кафедры)
\patitle{доцент, к.\,ф.-м.\,н.}
\paname{С.\,В.\,Миронов}

% Руководитель НИР
\nirtitle{доцент, к.\,п.\,н.} % степень, звание
\nirname{В.\,А.\,Векслер}

% Семестр (только для практики, для остальных типов работ не используется)
\term{2}

% Наименование практики (только для практики, для остальных типов работ не
% используется)
\practtype{учебная}

% Продолжительность практики (количество недель) (только для практики, для
% остальных типов работ не используется)
\duration{2}

% Даты начала и окончания практики (только для практики, для остальных типов
% работ не используется)
\practStart{01.07.2024}
\practFinish{13.01.2024}

% Год выполнения отчета
\date{2026}

\maketitle

% Включение нумерации рисунков, формул и таблиц по разделам (по умолчанию -
% нумерация сквозная) (допускается оба вида нумерации)
\secNumbering

\tableofcontents

% Раздел "Обозначения и сокращения". Может отсутствовать в работе
% \abbreviations
% \begin{description}
%     \item ... "--- ...
%     \item ... "--- ...
% \end{description}

% Раздел "Определения". Может отсутствовать в работе
% \definitions

% Раздел "Определения, обозначения и сокращения". Может отсутствовать в работе.
% Если присутствует, то заменяет собой разделы "Обозначения и сокращения" и
% "Определения"
% \defabbr

\intro

В настоящее время наблюдается стремительная цифровая трансформация сферы высшего образования, выраженная в повсеместном внедрении и использовании цифровых платформ для ведения образовательного процесса, в том числе систем управления обучением (LMS). Современные академические платформы эволюционировали от простых хранилищ методических материалов до комплексных распределенных систем, обеспечивающих интерактивное взаимодействие между студентами и преподавательским составом.

Ключевым элементом подобных платформ является подсистема конструирования и редактирования учебных курсов. При этом встречаются случаи, когда контент одного курса (структура модулей, текстовые блоки, задания, тесты) может администрироваться одновременно несколькими сотрудниками кафедры. В условиях многопользовательского доступа к распределенной веб-системе возникает важная инженерная проблема "--- обеспечение транзакционной целостности данных и предотвращение конфликтов при параллельном изменении одних и тех же сущностей. Без внедрения специализированных механизмов координации сессий редактирования может возникнуть ситуация коллизии данных или состояние гонки (race condition), когда изменения, внесенные одним пользователем, бесследно затираются запросами другого.

Дополнительная сложность заключается в специфике пользовательского опыта (UX) в веб-приложениях: стандартные механизмы транзакций в базах данных не подходят для того, чтобы изолировать процесс ручного редактирования контента, который может длиться часами. При этом необходимо обеспечить возможность версионирования учебных материалов "--- сохранения состояний курса в разные моменты времени с возможностью отката контента к предыдущим состояниям. Таким образом, разработка эффективного механизма редактирования курсов, сочетающего в себе координацию интерфейсов, строгую изоляцию данных на уровне СУБД и систему сохранения состояний, является актуальной научно-практической задачей в области программной инженерии.

Целью этой работы является проектирование и программная реализация отказоустойчивого и эффективного механизма редактирования учебных курсов с поддержкой изолированных черновиков и возможностью отката изменений для образовательной платформы <<MergeMinds>>.

Для достижения поставленной цели необходимо решить следующие задачи:

\begin{enumerate}
  \item изучить и систематизировать существующие теоретические подходы к управлению параллельным доступом к данным в многопользовательских системах (оптимистическое, пессимистическое и распределенное управление конкурентностью);
  \item проанализировать методы версионирования и хранения истории изменений сложных структур данных в реляционных системах управления базами данных;
  \item провести сравнительный анализ альтернативных архитектурных решений, выявить их достоинства, недостатки и компромиссы применительно к специфике образовательных платформ;
  \item разработать схему данных, алгоритмы взаимодействия компонентов (бэкенд-интерфейса и фронтенд-клиента) и протокол поддержания сессий редактирования на основе выбранного подхода;
  \item реализовать спроектированный механизм в рамках существующей архитектуры бэкенда образовательной платформы <<MergeMinds>>;
  \item оценить эффективность разработанного решения и его устойчивость к сетевым сбоям и нештатным сценариям поведения пользователей.
\end{enumerate}

% После введения — серии \section, \subsection и т.д.

\conclusion

% Отобразить все источники. Даже те, на которые нет ссылок.
% \nocite{*}

\bibliographystyle{ugost2003}
\bibliography{thesis}

% Окончание основного документа и начало приложений Каждая последующая секция
% документа будет являться приложением
\appendix

\end{document}
